% !TeX program = xelatex
\documentclass[UTF8,12pt,a4paper,fontset=fandol]{ctexart}
\usepackage[margin=25mm]{geometry}
\usepackage{amsmath,amssymb}
\usepackage[unicode,colorlinks=true,linkcolor=black,urlcolor=blue]{hyperref}
\numberwithin{equation}{section}
\allowdisplaybreaks[2]
\emergencystretch=2em
\title{公开视频课程炼化：最小讲义示例}
\author{整理者原创示例}
\date{}
\begin{document}
\maketitle
\section*{范围与核验说明}
此文用于检验中文排版与公式编号，不是任何视频课程的已完成整理。
原课事实须附来源；整理者解释须标明归属。编译成功不能代替数学与视觉审阅。
\tableofcontents
\section{换元为何有效}
\subsection{题目与第一直觉}
计算下式。根号内部的导数与分子中的一部分匹配，先拆出微分因子。
\begin{equation}\label{eq:question}
I=\int\frac{x^3}{\sqrt{1+x^2}}\,dx,\qquad x\in\mathbb R.
\end{equation}
\subsection{整理者解释与必要推导}
取新变量后，剩下的平方也能由新变量表达；这使函数和微分同时简化。
\begin{align}
u&=1+x^2, & du&=2x\,dx,\label{eq:substitution}\\
x^3\,dx&=\frac12(u-1)\,du,\\
I&=\frac12\int\left(u^{1/2}-u^{-1/2}\right)\,du\\
 &=\frac13u^{3/2}-u^{1/2}+C\\
 &=\frac13(1+x^2)^{3/2}-\sqrt{1+x^2}+C.\label{eq:answer}
\end{align}
\subsection{条件、误区与简单路线比较}
实数范围内根号有定义。注意式\eqref{eq:substitution}中的系数。
三角代换需要额外恒等式与回代，此处直接换元步骤更少。
无需强求全实轴上的单值反函数，所得结果可以直接求导核验。
\subsection{独立核验}
对式\eqref{eq:answer}的原函数求导，有
\begin{align}
F'(x)&=x\sqrt{1+x^2}-\frac{x}{\sqrt{1+x^2}}\\
&=\frac{x^3}{\sqrt{1+x^2}}.
\end{align}
这项解析核验与实际运行CAS分别记录。真实课程还须检查原题和来源。
\end{document}
